\documentclass[a4paper,12pt,oneside]{article}
\makeatletter
\def\input@path{{../../}}
\makeatother
\usepackage{discrete-practice}

\usepackage{booktabs,float}

\begin{document}

\PracticeHeader{1}{11}{Кодирование}

\Task[Коды Хаффмана и Шеннона--Фано]{Задание 1}
Имеется следующий алфавит и вероятностная модель:

\begin{table}[H]
\centering
\begin{tabular}{cc}
\toprule
Символ & Вероятность \\
\midrule
a & 0.35 \\
b & 0.17 \\
c & 0.17 \\
d & 0.16 \\
e & 0.15 \\
\bottomrule
\end{tabular}
\end{table}

\subsection*{a) Код Шеннона--Фано}
Найти код Шеннона--Фано. Определить среднюю длину кодового слова.

\subsection*{b) Код Хаффмана}
Построить код Хаффмана. Определить среднюю длину кодового слова.

\subsection*{c) Сравнение кодов}
Определить, у какого символа код Шеннона--Фано длиннее, чем код Хаффмана, и объяснить, почему выгоднее выбрать более короткое кодовое слово.

\Task[Код Шеннона--Фано]{Задание 2}
Построить код для системы из семи букв $\{A,B,C,D,E,F,G\}$ с вероятностями
$0.1$, $0.2$, $0.05$, $0.3$, $0.05$, $0.15$, $0.15$.

Найти среднее количество разрядов. Декодировать последовательность:
\[
10011101001000111101110101111000.
\]

\Task[Код Хаффмана для сочетаний]{Задание 3}
Алфавит из трёх символов: $0.75$, $0.1$, $0.15$. Построить:
\begin{alphlist}
    \item код для отдельных букв;
    \item код для двухбуквенных сочетаний.
\end{alphlist}
Определить средние длины и коэффициенты оптимальности.

\Task[Код Хаффмана для блоков]{Задание 4}
Алфавит $\{A,B,C\}$ с вероятностями $p(A)=0.6$, $p(B)=0.3$, $p(C)=0.1$.
Построить коды при кодировании:
\begin{alphlist}
    \item по одной букве;
    \item по две буквы;
    \item по три буквы.
\end{alphlist}
Сравнить эффективность.

\Task[Арифметическое кодирование]{Задание 5}
Получено сообщение:
\[
0110100110101.
\]
Известно, что в исходном сообщении было: $3$ A, $1$ C, $1$ D, $1$ Y, $2$ N. Декодировать.

\end{document}
